\VerdierPhysicalPageAnchor{0036}%
\VerdierSemanticAnchor{ch1-1-1-2-degree-e-morphism-cont}
de degré $e$ vérifiant les propriétés de compatibilité usuelles.

\VerdierSemanticAnchor{ch1-1-1-3}
\noindent\textbf{Définition 1.1.3.}\quad
Une catégorie graduée de type $G$ est dite \emph{additive} si la catégorie
fibre $\mathcal C_e$ est additive [20].
\VerdierSemanticAnchor{ch1-1-1-3-additive-functor}
Un foncteur gradué entre catégories graduées additives est dit
\emph{additif} si le foncteur restreint aux fibres est additif.

\VerdierSemanticAnchor{ch1-1-1-3-cadd}
Soit $\mathcal C$ une catégorie graduée de type $G$ additive. La catégorie
$\mathcal C$ dépouillée de sa structure fibrée \emph{n'est pas additive}
lorsque $G\neq\{e\}$. On peut toutefois associer à $\mathcal C$ une
catégorie additive en procédant comme suit : pour tout couple $X,Y$
d'objets de $\mathcal C$ et tout élément $g$ de $G$, on vérifie que
l'ensemble $\operatorname{Hom}_{\mathcal C}^{g}(X,Y)$ est muni
naturellement d'une structure de groupe commutatif. Définissons alors une
nouvelle catégorie $\mathcal C_{\mathrm{add}}$ en prenant comme ensemble
d'objets l'ensemble des objets de $\mathcal C$, et comme ensemble de
morphismes entre deux objets $X$ et $Y$ :

\VerdierSemanticAnchor{ch1-1-1-3-eq1}
\begingroup
\makeatletter
\tagsleft@true
\let\veqno\@@leqno
\makeatother
\begin{equation}
  \operatorname{Hom}_{\mathcal C_{\mathrm{add}}}(X,Y)
  =\bigoplus_{g\in G}\operatorname{Hom}_{\mathcal C}^{g}(X,Y).
  \tag{1.1.3.1}
\end{equation}
\endgroup

\VerdierSemanticAnchor{ch1-1-1-3-cadd-composition}
La composition des morphismes dans $\mathcal C_{\mathrm{add}}$ se définit
de la manière évidente et on vérifie que $\mathcal C_{\mathrm{add}}$ est
une catégorie additive.
\VerdierSemanticAnchor{ch1-1-1-3-fadd}
Étant donné de même un foncteur cartésien additif
$F\colon\mathcal C\to\mathcal C'$ entre deux catégories graduées
additives, on construit par la même méthode un foncteur additif
$F_{\mathrm{add}}\colon\mathcal C_{\mathrm{add}}\to
\mathcal C'_{\mathrm{add}}$. Nous n'utiliserons jamais cette construction.

\VerdierSemanticAnchor{ch1-1-1-3-restrictive}
Remarquons que la définition des catégories graduées adoptée ici est plus
restrictive que celle qui s'impose naturellement à l'esprit. En plus des
structures évidentes (1.1.2.1) et (1.1.2.2), nous imposons la propriété
(1.1.2.3).

\VerdierSemanticAnchor{ch1-1-2}
\subsection{Opération à isomorphisme près d'un groupe sur une catégorie}

\VerdierSemanticAnchor{ch1-1-2-1}
\noindent\textbf{1.2.1.}\quad
Soit $\mathcal C$ une catégorie graduée de type $G$. Choisissons un clivage
normalisé de la catégorie fibrée
$\mathcal C\xrightarrow{\deg}\mathcal G$ [18], \emph{i.e.} pour tout
$g\in G$ choisissons un foncteur changement de base :

\VerdierSemanticAnchor{ch1-1-2-1-eq1}
\begingroup
\makeatletter
\tagsleft@true
\let\veqno\@@leqno
\makeatother
\begin{equation}
  \begin{cases}
    T(g)\colon\mathcal C_e\longrightarrow\mathcal C_e,\\[0.4em]
    T(e)=\text{identité de }\mathcal C_e.
  \end{cases}
  \tag{1.2.1.1}
\end{equation}
\endgroup

\VerdierSemanticAnchor{ch1-1-2-1-transition}
Pour tout couple $(g,h)$ d'éléments de $G$, on a [18] un isomorphisme de
foncteurs :

\VerdierSemanticAnchor{ch1-1-2-1-eq2}
\begingroup
\makeatletter
\tagsleft@true
\let\veqno\@@leqno
\makeatother
\begin{equation}
  \begin{cases}
    c(h,g)\colon T(hg)\longrightarrow T(g)T(h),\\[0.4em]
    c(g,e)=c(e,g)=\text{identité de }T(g)
  \end{cases}
  \tag{1.2.1.2}
\end{equation}
\endgroup
